\section{Conclusion}

<<<<<<< HEAD
This study demonstrates that a consumer-grade, wireless 3-user Meta Quest 3 setup can meet selected technical benchmarks relevant to co-located training in a controlled testbed.

\subsection{Key Findings}

The technical validation (RQ5) confirmed that the 3-user co-located system met the evidence-based benchmarks used as reference targets in our testbed at the \emph{mean} level. Network latency averaged $58.2 \pm 21.3$ms, remaining below Van Damme et al.'s RTT target~\cite{VanDammeSam2024Iolo} for good collaborative quality of experience, although tail latency exceeded the 75ms target at P95. Under stationary baseline conditions, calibration error remained well below Reimer et al.'s $<10$mm reference threshold~\cite{ReimerDennis2021CfSV} (mean $1.21 \pm 0.90$mm; max 5.78mm on the one headset that recorded non-zero values). Frame rate remained stable around the selected 72Hz target ($72.2$fps) throughout the extended run. Regarding long-term stability (RQ5), thermal analysis showed substantial warming (max $52.6^{\circ}$C) while performance remained stable.

\subsection{Practical Implications}

The system demonstrates significant cost-effectiveness. At approximately \$2,000 in hardware cost for three Meta Quest 3 headsets plus consumer networking equipment, it enables a portable testbed for co-located collaborative VR without dedicated VR-room infrastructure. Wireless operation offers deployment flexibility by eliminating cable management complexity, enabling training in authentic physical environments rather than dedicated VR rooms. Additionally, the measured calibration error under stationary baseline conditions indicates substantial safety margin relative to the 10mm threshold used as a reference.
=======
This study validates consumer-grade wireless VR headsets for professional co-located training applications, demonstrating that Meta Quest 3 systems achieve enterprise-grade performance benchmarks at 15--20\% of traditional costs.

\subsection{Key Findings}

The technical validation (RQ1) confirmed that the 3-user co-located system met the evidence-based requirements used as reference targets in our testbed. Network latency averaged $58.2 \pm 21.3$ms, remaining within Van Damme et al.'s threshold~\cite{VanDammeSam2024Iolo} for good Quality of Experience. Under stationary baseline conditions, calibration error remained well below Reimer et al.'s $<10$mm reference threshold~\cite{ReimerDennis2021CfSV} (mean $1.21 \pm 0.90$mm; max 5.78mm on the headset reporting non-zero calibration error). Furthermore, frame rates remained stable around the native 72Hz target ($72.2$fps) throughout the extended run. Regarding long-term stability (RQ4), thermal analysis revealed substantial warming (max $52.6^{\circ}$C) while performance remained stable.

\subsection{Practical Implications}

The system demonstrates significant cost-effectiveness. At approximately \$1,500 total hardware cost (3$\times$ Quest 3 headsets) plus consumer networking equipment, it delivers 98\% of enterprise VR room capabilities at just 15--20\% of typical costs, democratizing access to high-fidelity collaborative VR training. Wireless operation offers deployment flexibility by eliminating cable management complexity, enabling training in authentic physical environments rather than dedicated VR rooms. Additionally, achieving single-digit millimeter calibration error under stationary baseline conditions provides substantial safety margins for collision prevention in co-located training scenarios.
>>>>>>> 32df077fa8bbb82b5374f4992fe0114abd8167f7

\subsection{Limitations}

While mean latency was acceptable, occasional spikes indicate the need for robust error handling in production environments to ensure network stability. Additionally, this study validated 3-user configurations; scaling to 4 or more users may introduce network congestion, potentially requiring additional infrastructure to maintain performance.
<<<<<<< HEAD
Because the drift evaluation (RQ2) was conducted with stationary headsets, the reported millimeter-scale alignment should be treated as a best-case baseline and does not, by itself, guarantee collision safety under active user motion.


\subsection{Future Work}

\subsubsection{Scalability and Infrastructure}
Future research should systematically evaluate 4-, 5-, and 6-user configurations to determine practical upper limits for scalability. In addition to networking load, these studies should document the infrastructure assumptions required for repeatable low-latency local sessions, such as channel selection, and contention from other nearby devices.

\subsubsection{Comparative Training Evaluation}
To move from technical feasibility to training relevance, controlled studies should compare training outcomes between co-located VR, remote VR, and traditional instruction. For collaborative medical training in particular, structured pilot studies with medical students and clinical staff in short, repeatable scenarios can test whether co-location improves communication, role clarity, and shared situational awareness.

\subsubsection{Teamwork, Safety, and Motion-Robustness}
Future prototypes should incorporate evaluation measures aligned with teamwork and safety objectives, including task completion time, error rates, observer-assessed teamwork/communication rubrics, subjective workload, and explicit safety measures such as near-miss events and minimum spacing between users. Importantly, calibration drift should be re-evaluated under active motion to quantify how alignment error impacts both perceived realism and collision risk.

\subsubsection{Physical Realism and MR Integration}
Purely virtual training remains limited by reduced sensory fidelity. A promising direction is to combine co-located VR with physical props, such as real tools, reference surfaces, or simple mannequins, so that key actions are grounded in touch, while retaining the flexibility and repeatability of virtual content.

\subsubsection{Spatial Awareness and Coordination Aids}
Finally, future work should explore spatial awareness and coordination aids tailored to co-located clinical workflows. This includes testing external-view interfaces, such as world-in-miniature or map-style views for role-based coordination, as well as procedure-specific cues that help users anticipate teammates' movements in tight spaces. These additions should be evaluated not only for usability, but also for whether they reduce collision risk without distracting from the primary task.


\subsection{Concluding Remarks}

This research demonstrates that current-generation VR hardware has matured to the point of viability for professional training applications. By systematically validating performance against evidence-based benchmarks, we provide a methodological framework for future co-located multi-user VR system development. The combination of affordable standalone headsets, robust wireless networking, and sophisticated spatial tracking enables a new practical foundation for portable, collaborative VR training.
=======

\subsection{Future Work}

Future research should systematically evaluate 4-, 5-, and 6-user configurations to determine practical upper limits for scalability. It is also crucial to conduct controlled studies comparing training outcomes between co-located VR, remote VR, and traditional methods to assess learning effectiveness.

\subsection{Concluding Remarks}

This research demonstrates that current-generation consumer VR hardware has matured to the point of viability for professional training applications. By systematically validating performance against evidence-based benchmarks, we provide a methodological framework for future co-located multi-user VR system development. The convergence of affordable standalone headsets, robust wireless networking, and sophisticated spatial tracking enables a new paradigm of portable, collaborative VR training.

>>>>>>> 32df077fa8bbb82b5374f4992fe0114abd8167f7
